\chapter{NextGenPB}
\label{sec:nextgenpb}

At the time of writing, Next Generation Poisson-Boltzmann (NextGenPB or NGPB for short) \cite{diflorio2025nextgenpb} is the most recent innovation in the field of LPBE solvers, aiming to improve accuracy while maximizing computational efficiency. 
It is part of the family of grid- or FEM-based solvers, featuring local (de)refinement.
It is purposefully modular and built upon free and open libraries, making improvements to individual steps easier, though, as I will explain later, this has to be taken with a grain of salt.
Its main attractions are the use of NanoShaper and an enhanced analytical approach, allowing for much more accurate energy calculation, beating out many older LPBE solvers in both accuracy and speed --- an impressive feat that has gained the attention of biomolecular chemists at our technical university, who seek to replace their T-ABPS with NGPB.
% yes I actually put that em dash there

%Just like other LPBE solvers, 
To avoid the self-energy singularity, NGPB computes the electrostatic energies individually, where each of them is never evaluated at the atom's center.
This is an entirely new approach, which avoids reformulating, doing additional differencing or using a different model other than point charges.
Thus the total electrostatic energy NGPB computes is a sum of the Coulomb, polarization and ionic energy.

The Coulomb energy is the typical interaction between two charged particles $q_i$ and $q_j$ without self-interaction, so $\lVert \mathbf{r}_i - \mathbf{r}_j \rVert$ in Equation~\eqref{eq:Coulomb-energy} never equals 0. 

\begin{equation}
    E_{coul} = \sum_{i=1}^{N_{atoms}} \sum_{j<i} \frac{q_i q_j}{4\pi\varepsilon_m \lVert \mathbf{r}_i - \mathbf{r}_j \rVert}
    \label{eq:Coulomb-energy}
\end{equation}

The polarization energy arises from polarization charges at the surface where the dielectric constant jumps. It is always evaluated on the intersection point $r_p$ between the surface and the grid.

\begin{equation}
    E_{pol} = \frac{1}{2} \sum_{i=1}^{N_{atoms}} \sum_{p=1}^{N_{Ic}} \frac{q_i q_p}{4\pi\varepsilon_0 \lVert \mathbf{r}_p - \mathbf{r}_i \rVert}
    \label{eq:polarization-energy}
\end{equation}

The ionic energy is generated by the interaction between the solvent counterions and the charges of the solute, and it is very small compared to the polarization and Coulomb energies.
Computing it would require doing an expensive volume integration over the entire solute volume, which is why it was either avoided, computed implicitly by evaluating the total energy directly, or computed indirectly through multiple runs with different ionic strengths.
However, using Green's identities, the authors of NGPB converted the volume integral into a surface integral over $\Gamma$.

\begin{equation}
    E_{ion} = \frac{1}{2} \sum_{i=1}^{N_{atoms}} q_i \left( \int_{\Gamma} \phi(\tilde{\mathbf{r}}) \frac{(\tilde{\mathbf{r}} - \mathbf{r}_i) \cdot \mathbf{n}(\tilde{\mathbf{r}})}{4\pi \lVert \tilde{\mathbf{r}} - \mathbf{r}_i \rVert^3} \, dS - \frac{1}{\varepsilon_s} \int_{\Gamma} \frac{\mathbf{D}(\tilde{\mathbf{r}}) \cdot \mathbf{n}(\tilde{\mathbf{r}})}{4\pi \lVert \tilde{\mathbf{r}} - \mathbf{r}_i \rVert} \, dS \right)
    \label{eq:ionic-energy}
\end{equation}

The total energy is then simply the sum of all three contributions:

\begin{equation}
    E_{tot} = E_{coul} + E_{pol} + E_{ion}
    \label{eq:total-energy}
\end{equation}

\section{Libraries}
\label{sec:nextgenpb-libraries}

\subsection{P4est}
\label{sec:nextgenpb-libraries-p4est}

P4est \cite{carsten_burstedde_2024_10839051} is a library for dynamic management of adaptive quadtrees, while octrees are taken care of by its submodule p8est.
It is designed for parallel work and scales well on many-core systems, though unfortunately, it has not been officially adapted to work on GPUs.

\subsection{NanoShaper}
\label{sec:nextgenpb-libraries-nanoshaper}

NanoShaper \cite{abate2025nanoshaperweb} is a surface reconstruction and pocket detection tool aimed at biomolecular systems.
It is compatible with patch-based, implicit surfaces or triangulated surfaces and works by casting rays in a grid pattern, computing both surface intersections and normals. 
The authors leverage this by aligning the grid of their problem space with that of NanoShaper.

\section{Submodules}
\label{sec:nextgenpb-submodules}

\subsection{Energy Calculation}
\label{sec:nextgenpb-submodules-energy-calculation}